% Introduction: motivation, VVUQ principles, thesis, contributions.
% Instruction only. The Introduction precedes the Table of Contents.

[PROCESSING INSTRUCTION (Introduction). Open with the clinical and
technical need: physical AI is entering oncology trials through in
silico simulation, autonomous code generation, and eventually robotic
procedures, where an error that ships is far costlier than one caught
before shipping. Then define the three VVUQ dimensions in the oncology
trial context, briefly and in plain terms, using exactly these framings:
Verification asks ``Did we build the code correctly?''; Validation asks
``Did we build the right model?''; and Uncertainty Quantification asks
``How confident are we in the results?''. Make clear that code
generation answers none of these on its own.

State the thesis next, in the author's voice: the LLM VVUQ process must
be more substantial than the code generation it gates, and being so is
what makes autonomous oncology trial deliverables faster, cheaper, and
more rigorous than conventional verification. Then give the three to
four contributions of this work, each mapped forward to a later section:
(i) a single repeatable pipeline of five established methods (Methods);
(ii) a VVUQ gate held to a higher standard than generation, exercised
across its full accept, block, and escalate surface (Results); (iii) an
honest account of what ran, what was approximated, and what could not run
(Limitations); and (iv) a Stage 2 2030 PDAC deployment reference
(Results, Discussion). Close by positioning this platform as the next
step beyond the prior humanoid instruction to code to execution to paper
workflow, which proved the workflow but left the assurance and per commit
automation to be strengthened here.

SOURCE FILES (Table~\ref{tab:intro_sources}; read first, then write in your
own words):

\begin{table}[H]
\caption{Introduction source map (abbreviated paths).}
\label{tab:intro_sources}
\centering
\begin{tabular}{L{4.6cm} L{8.8cm}}
\toprule
\textbf{Source (under repo root)} & \textbf{Use for} \\
\midrule
  vvuq/ : README + 4 modules & The three VVUQ dimensions and the gate principle, stated plainly. \\
  execution/README.md : Thesis section & The thesis wording and the assurance over generation argument. \\
  inputs/Paper-VVUQ-1, -2 : README + chunk 01 & Background on VVUQ for in silico trials; the 117 DOI input corpus. \\
  inputs/VVUQ-Research-1, -2 : README + chunk 1 & Additional VVUQ and digital twin framing for oncology. \\
  Clinical-AI-Demos 07-humanoid final-paper & The prior workflow this paper extends (cite, do not copy). \\
\bottomrule
\end{tabular}
\end{table}

CITATIONS: support the VVUQ definitions and the precision medicine
framing with \cite{sel2025vvuq}, \cite{zaid2025pbpk}, and
\cite{moradi2026twins}; the credibility framing with \cite{asmevv40} and
\cite{fda2023credibility}; the autonomous generation context with
\cite{trooskens2026compiled}, \cite{carlisle2026pipeline}, and
\cite{yan2026clinagent}; and the platform lineage with
\cite{kawchak2026sponsor}, \cite{kawchak2025qsppdac}, and
\cite{repo-cancer-automated}.

HOW TO PROCESS: keep the VVUQ definitions to a few sentences each; do not
restate Table~\ref{tab:intro_sources}, reference it. Connect the thesis
here to the Results and Discussion so the paper reads as one argument.]
